% QCE 2026 poster submission -- PAGE 1 (sections 1-4): title, authors+contact,
% 200-250 word abstract (for website), relevance to QCE topics.
% Standalone 1-page cover; prepend to main.tex (2pp) + poster.pdf for the package.
\documentclass[10pt, conference]{IEEEtran}
\IEEEoverridecommandlockouts
\usepackage{amsmath,amssymb}
\usepackage{xcolor}

\begin{document}

% --- 1. Poster title ---
\title{Difficulty Analysis of Posiform-Planted QUBO Benchmarks\\
via Energy and Hamming Metrics}

% --- 2. Authors + representative contact ---
\author{
\IEEEauthorblockN{Yideun Kim}
\IEEEauthorblockA{
Department of Computer Engineering\\
Kangwon National University\\
Chuncheon, Republic of Korea\\
Email: rladlems1031@kangwon.ac.kr}
\and
\IEEEauthorblockN{Dongjae Lee$^{\dagger}$}
\IEEEauthorblockA{
Department of Convergence Security\\
Kangwon National University\\
Chuncheon, Republic of Korea\\
Email: dongjae.lee@kangwon.ac.kr}
}
\maketitle
\renewcommand{\thefootnote}{}
\footnotetext{$^{\dagger}$Corresponding author and representative contact:
Dongjae Lee (\texttt{dongjae.lee@kangwon.ac.kr}).}
\renewcommand{\thefootnote}{\arabic{footnote}}

% --- 3. Abstract (200-250 words, for website posting if accepted) ---
\begin{abstract}
Benchmarking quantum annealers and classical heuristics needs problem instances
whose optimal solution is known in advance but that stay hard to solve. Posiform
planting generates Quadratic Unconstrained Binary Optimization (QUBO) instances
with one unique planted optimum. A hardened variant raises their difficulty. It
adds random spin-glass terms and scales the planted component by a coefficient
$\alpha$. A practical obstacle is measurement. Hardness is usually graded by the
ground-state success rate. But across the whole hard regime this rate is $0\%$,
so it cannot distinguish instances, solvers, or values of $\alpha$. This poster
presents a measurement protocol for the zero-success regime. Instead of asking
whether a run succeeded, we ask how close the failing runs come to the planted
optimum. We use two quantities, computed over failures only and on a fixed set
of the hardest instances. The first is the energy gap to the optimum. The second
is the Hamming distance from the planted answer. Together they turn the flat
$0\%$ rate into a continuous, rankable difficulty curve. On a D-Wave Pegasus
graph ($n=5612$ qubits), the curve is non-monotone in $\alpha$. It rises to a
peak and then declines. A single supporting run on a D-Wave Advantage annealer
reproduces the same shape, more weakly. This is consistent with the structure
reflecting the instances rather than simulated annealing alone. The protocol
gives a finer way to grade planted QUBO benchmarks.
\end{abstract}

\begin{IEEEkeywords}
QUBO benchmarking, quantum annealing, simulated annealing, posiform planting,
performance evaluation, NISQ, quantum optimization.
\end{IEEEkeywords}

% --- 4. Relevance to QCE topics ---
\section*{Poster Relevance to IEEE Quantum Week (QCE 2026)}
This work addresses the benchmarking, performance evaluation, and verification
of programmable quantum annealers. These are core themes of IEEE Quantum Week.
Its subject is the QUBO problem, the native input format of D-Wave quantum
annealers and a central model for quantum and quantum-inspired optimization. The
proposed failure-conditioned metrics target the NISQ-era regime, where solvers
seldom reach the true ground state and a binary success rate is uninformative.
They still rank instances, solvers, and a hardware-relevant difficulty
parameter. The study spans both classical (simulated annealing) and quantum
(D-Wave Advantage) solvers under a common measure, which supports cross-platform
comparison. It maps directly to two QCE topics: \emph{adiabatic quantum computation \&
quantum annealing} (Quantum Computing), the model that D-Wave devices implement;
and \emph{benchmarks} (Quantum Systems Software), since the failure-conditioned
protocol and open instance set form a benchmark for the zero-success regime. All
code and benchmark instances are released openly
(\texttt{github.com/YIDEUNKIM/qubo\_dataset}) to support reproducibility.

\end{document}
